\section{表示把群变成矩阵群}
\label{sec:04-Discrete-Mathematics/09-Representation-Theory/01}

桌上有一块正三角形瓷砖。把它拿起来再放回原位，可行的姿势有六种：原样放回，绕中心转 \(120^\circ\) 或 \(240^\circ\)，沿三条对称轴各翻一次。这六个动作在复合下封闭，构成一个六元素的群，同构于 \(S_3\)。\hyperref[ch:032]{``群与对称：可逆操作怎样组织不变性''}已经把这类对象认过一遍。

把乘法表写全，是六行六列的方格，每格填一个动作的名字。然后呢？想问``这个群内部有没有几套彼此不干扰的机制''，或者``手上这两个群是不是同一个东西换了记号''，乘法表帮不上忙。格子里装的是名字：名字不能相加，没有特征值，也谈不上远近。

出路是给每个动作配一个矩阵。转 \(120^\circ\) 在平面上就是旋转矩阵，沿轴翻转就是反射矩阵，两个动作先后做对应两个矩阵相乘。走完这一步，群从一张组合表格变成一族线性变换，迹、特征值、不变子空间、直和分解——线性代数攒下的整套工具立刻可用。本单元后面的全部内容，都是这一步换来的。

\subsection{唯一的要求是把乘法搬过去}

\begin{definition}
设 \(V\) 是复数域上的有限维向量空间。群 \(G\) 在 \(V\) 上的表示指一个映射 \(\rho:G\to GL(V)\)，满足
\[
\rho(gh)=\rho(g)\rho(h)
\qquad\text{对所有 }g,h\in G,
\]
其中 \(GL(V)\) 是 \(V\) 上全体可逆线性变换构成的群。\(\dim V\) 称为该表示的维数。
\end{definition}

条件只有这一条，别的都是它的推论：在等式里取 \(h=e\) 得 \(\rho(e)=I\)，再取 \(h=g^{-1}\) 得 \(\rho(g^{-1})=\rho(g)^{-1}\)。所以不必额外要求``单位元映到单位矩阵''，同态自己会安排。

\begin{center}
\begin{tikzpicture}[x=1cm,y=1cm,font=\small,>=stealth]
  \draw[black!45,fill=blue!7,rounded corners=4pt] (-0.3,-0.3) rectangle (3.6,5.0);
  \node[black!70] at (1.65,4.6) {群 \(G\)：动作};
  \node at (1.65,3.6) {\(r\)：转 \(120^\circ\)};
  \node at (1.65,2.2) {\(s\)：沿中线翻};
  \node at (1.65,0.7) {\(rs\)：先翻，再转};
  \draw[->,black!60] (1.65,1.75) -- (1.65,1.15);
  \node[right,black!65,font=\footnotesize] at (1.75,1.45) {复合};

  \draw[black!45,fill=orange!8,rounded corners=4pt] (5.6,-0.3) rectangle (11.4,5.0);
  \node[black!70] at (8.5,4.6) {\(GL_2(\C)\)：可逆矩阵};
  \node at (8.5,3.6) {\(\rho(r)=\begin{pmatrix}-\tfrac12&-\tfrac{\sqrt3}{2}\\[2pt]\tfrac{\sqrt3}{2}&-\tfrac12\end{pmatrix}\)};
  \node at (8.5,2.2) {\(\rho(s)=\begin{pmatrix}1&0\\0&-1\end{pmatrix}\)};
  \node at (8.5,0.7) {\(\rho(rs)=\rho(r)\,\rho(s)\)};
  \draw[->,black!60] (8.5,1.6) -- (8.5,1.1);
  \node[right,black!65,font=\footnotesize] at (8.6,1.35) {矩阵乘法};

  \draw[->,thick] (3.7,3.6) -- (5.5,3.6) node[midway,above,font=\footnotesize] {\(\rho\)};
  \draw[->,thick] (3.7,2.2) -- (5.5,2.2) node[midway,above,font=\footnotesize] {\(\rho\)};
  \draw[->,thick] (3.7,0.7) -- (5.5,0.7) node[midway,above,font=\footnotesize] {\(\rho\)};
\end{tikzpicture}

\emph{左右两侧各有一套``先做这个再做那个''的运算，\(\rho\) 要求两侧对得上：先在左边复合再翻译，与先翻译再在右边相乘，结果必须一样。这张图后面会反复用到——判断一个矩阵族是不是表示，看的永远是这个方框合不合拢。}
\end{center}

\subsection{四个例子，先摆出来}

平凡表示把每个 \(g\) 都送到 \(1\times1\) 的单位矩阵。它丢掉了群的全部信息，但每个群都有，后面做正交投影时它是必备的一根坐标轴。

置换表示来自群作用。\(S_3\) 作用在三个坐标轴上，取 \(V=\C^3\)、让 \(\rho(g)\) 把基向量 \(e_i\) 送到 \(e_{g(i)}\)，得到三个坐标的重排：
\[
\rho\bigl((12)\bigr)=\begin{pmatrix}0&1&0\\1&0&0\\0&0&1\end{pmatrix},\qquad
\rho\bigl((123)\bigr)=\begin{pmatrix}0&0&1\\1&0&0\\0&1&0\end{pmatrix}.
\]

正则表示更奢侈：直接拿群元素当基，\(V=\C[G]\) 的维数是 \(|G|\)，\(\rho(g)\) 把基向量 \(h\) 送到 \(gh\)。它把整个群装进了一个向量空间，第三节会看到它同时装下了这个群的每一个不可约表示。

上面那张图里的二维表示不属于以上任何一类。它不重排基向量，而是把整个平面转起来、翻过去：坐标被混合了。可以核对 \(\rho(r)^3=I\)、\(\rho(s)^2=I\) 以及 \(\rho(s)\rho(r)\rho(s)=\rho(r)^{-1}\)，与三角形上那三个动作的关系完全吻合。这正是表示比群作用宽的地方——群作用只能置换一堆点，表示允许任意的线性混合。二维旋转群 \(SO(2)\) 在 \(\R^2\) 上的自然表示同理，没有哪组基向量是被``排列''的。

\subsection{换基不改变表示}

同一个线性算子换一组基就换一副矩阵面孔，表示也一样。若存在可逆的 \(T:V\to V'\) 使
\[
\rho'(g)=T\rho(g)T^{-1}
\qquad\text{对所有 }g\in G,
\]
就说 \(\rho\) 与 \(\rho'\) 等价。注意 \(T\) 只有一个，要同时替所有群元素完成换装。这与\hyperref[sec:03-Linear-Algebra/05-Eigenvalues-and-Eigenvectors/02]{``对角化与动力系统：把耦合演化变成独立缩放''}里的相似变换是同一件事，只是那里换基是为了让一个矩阵变简单，这里要让一整族矩阵同时变简单。

于是有意义的研究对象是等价类，不是某一组具体矩阵。凡是随基改变的量都不配当表示的性质；反过来，凡是在相似变换下不变的量都值得留意。迹就是这样一个量，第三节整篇都在用它。

\subsection{矩阵能看见乘法表看不见的东西}

子空间 \(W\subseteq V\) 叫做不变子空间，如果 \(\rho(g)W\subseteq W\) 对每个 \(g\) 成立。回到 \(S_3\) 的置换表示：向量 \((1,1,1)^{\trans}\) 对三个坐标一视同仁，任何重排都动不了它，所以它张成的直线是一维不变子空间。这样的直线在乘法表上根本不存在——它是矩阵语言带来的新观察。

有了不变子空间就能拆。挑一组基，先排满 \(W\) 再补齐其余，所有 \(\rho(g)\) 同时变成分块上三角的形状：左上块记录 \(W\) 上发生的事，右下块记录商空间上发生的事，左下角恒为零。若 \(V\) 除了 \(\{0\}\) 和自身之外再无不变子空间，就说 \(\rho\) 不可约，拆不动了。

一维表示自动不可约。正则表示在 \(|G|>1\) 时一定可约，因为所有坐标之和相等的那条直线被群平移固定，平凡表示藏在里面。

\subsection{底域取复数不是随手一挑}

本书默认 \(V\) 建在 \(\C\) 上，两个理由都很实在。复数域代数闭，任何 \(\rho(g)\) 都有特征值，找不变方向时不会因为特征多项式无根而卡住；而且对有限群，\(|G|\) 在 \(\C\) 中可逆，下一节的平均化构造才跑得通。

换掉底域，图景立刻变形。实数域上平面旋转没有实特征向量，二维的实不可约表示比比皆是；有限域上若域的特征整除 \(|G|\)，连``可以拆干净''都不再成立。复数域上的理论是最干净的那一层，也是理解其余各层的参照系。

\subsection{表示太多，需要原子}

同一个群的表示要多少有多少：维数随便取，还能把小表示塞进大空间里用零填满四周。这种自由度让``研究一个群的全部表示''听起来像空话。要让它成为可做的事，得先回答一个问题——这些表示里有没有一批拆不动的基本款，其余的都由基本款拼出来？

这个问题的答案对下游影响很大。等变神经网络的一层要与群作用交换，它的输入输出各自带着一个表示；能不能把这层的参数写清楚，取决于这两个表示各由哪些基本款组成。下一节从刚才那个三维例子接着拆，看看拆到什么地方停下来，以及凭什么保证一定拆得干净。
